% problem_id: S001-duality-lemma-shriek-well-defined
% source_id: S001 (Stacks Project)
% source_file: duality.tex
% statement_locator: duality.tex lines 3768--3773
% proof_locator: duality.tex lines 3775--3941
% env: lemma
% id_note: label 跨文件重名 ⇒ id 用文件限定形式
% pairing: tight (verified)
% status: complete (statement + author proof verbatim + reason + locators)
%
\begin{lemma}
\label{lemma-shriek-well-defined}
In Situation \ref{situation-shriek} let $f : X \to Y$ be a morphism
of $\textit{FTS}_S$. The functor $f^!$ is, up to canonical isomorphism,
independent of the choice of the compactification.
\end{lemma}

\begin{proof}
The category of compactifications of $X$ over $Y$ is defined
in More on Flatness, Section \ref{flat-section-compactify}.
By More on Flatness, Theorem \ref{flat-theorem-nagata} and
Lemma \ref{flat-lemma-compactifyable} it is nonempty.
To every choice of a compactification
$$
j : X \to \overline{X},\quad \overline{f} : \overline{X} \to Y
$$
the construction above associates the functor $j^* \circ \overline{a} :
D_\QCoh^+(\mathcal{O}_Y) \to D_\QCoh^+(\mathcal{O}_X)$
where $\overline{a}$ is the right adjoint of $R\overline{f}_*$
constructed in Lemma \ref{lemma-twisted-inverse-image}.

\medskip\noindent
Suppose given a morphism $g : \overline{X}_1 \to \overline{X}_2$
between compactifications $j_i : X \to \overline{X}_i$ over $Y$
such that $g^{-1}(j_2(X)) = j_1(X)$\footnote{This may fail
with our definition of compactification. See
More on Flatness, Section \ref{flat-section-compactify}.}.
Let $\overline{c}$ be the right adjoint of
Lemma \ref{lemma-twisted-inverse-image} for $g$.
Then $\overline{c} \circ \overline{a}_2 = \overline{a}_1$
because these functors are adjoint to
$R\overline{f}_{2, *} \circ Rg_* = R(\overline{f}_2 \circ g)_*$.
By (\ref{equation-sheafy}) we have a canonical transformation
$$
j_1^* \circ \overline{c} \longrightarrow j_2^*
$$
of functors
$D^+_\QCoh(\mathcal{O}_{\overline{X}_2}) \to D^+_\QCoh(\mathcal{O}_X)$
which is an isomorphism by Lemma \ref{lemma-proper-noetherian}.
The composition
$$
j_1^* \circ \overline{a}_1 \longrightarrow
j_1^* \circ \overline{c} \circ \overline{a}_2 \longrightarrow
j_2^* \circ \overline{a}_2
$$
is an isomorphism of functors which we will denote by $\alpha_g$.

\medskip\noindent
Consider two compactifications $j_i : X \to \overline{X}_i$, $i = 1, 2$
of $X$ over $Y$. By More on Flatness, Lemma
\ref{flat-lemma-compactifications-cofiltered} part (b)
we can find a compactification $j : X \to \overline{X}$
with dense image and morphisms
$g_i : \overline{X} \to \overline{X}_i$ of compactifications.
By More on Flatness, Lemma
\ref{flat-lemma-compactifications-cofiltered} part (c)
we have $g_i^{-1}(j_i(X)) = j(X)$. Hence we get isomorpisms
$$
\alpha_{g_i} :
j^* \circ \overline{a}
\longrightarrow
j_i^* \circ \overline{a}_i
$$
by the previous paragraph. We obtain an isomorphism
$$
\alpha_{g_2} \circ \alpha_{g_1}^{-1} :
j_1^* \circ \overline{a}_1 \to j_2^* \circ \overline{a}_2
$$
To finish the proof we have to show that these isomorphisms are well defined.
We claim it suffices to show the composition of isomorphisms constructed
in the previous paragraph is another (for a precise statement see the next
paragraph). We suggest the reader check this is true on a napkin, but we
will also completely spell it out in the rest of this paragraph.
Namely, consider a second choice of a compactification
$j' : X \to \overline{X}'$ with dense image
and morphisms of compactifications $g'_i : \overline{X}' \to \overline{X}_i$.
By More on Flatness, Lemma \ref{flat-lemma-compactifications-cofiltered}
we can find a compactification $j'' : X \to \overline{X}''$
with dense image and morphisms of compactifications
$h : \overline{X}'' \to \overline{X}$ and
$h' : \overline{X}'' \to \overline{X}'$. We may even assume
$g_1 \circ h = g'_1 \circ h'$ and $g_2 \circ h = g'_2 \circ h'$.
The result of the next paragraph gives
$$
\alpha_{g_i} \circ \alpha_h = \alpha_{g_i \circ h} =
\alpha_{g'_i \circ h'} = \alpha_{g'_i} \circ \alpha_{h'}
$$
for $i = 1, 2$. Since these are all isomorphisms of functors
we conclude that $\alpha_{g_2} \circ \alpha_{g_1}^{-1} =
\alpha_{g'_2} \circ \alpha_{g'_1}^{-1}$ as desired.

\medskip\noindent
Suppose given compactifications $j_i : X \to \overline{X}_i$
for $i = 1, 2, 3$.  Suppose given morphisms
$g : \overline{X}_1 \to \overline{X}_2$ and
$h : \overline{X}_2 \to \overline{X}_3$ of compactifications
such that $g^{-1}(j_2(X)) = j_1(X)$ and $h^{-1}(j_2(X)) = j_3(X)$.
Let $\overline{a}_i$ be as above. The claim above means that
$$
\alpha_g \circ \alpha_h = \alpha_{g \circ h} :
j_1^* \circ \overline{a}_1 \to j_3^* \circ \overline{a}_3
$$
Let $\overline{c}$, resp.\ $\overline{d}$ be the right adjoint of
Lemma \ref{lemma-twisted-inverse-image} for $g$, resp.\ $h$.
Then $\overline{c} \circ \overline{a}_2 = \overline{a}_1$ and
$\overline{d} \circ \overline{a}_3 = \overline{a}_2$
and there are canonical transformations
$$
j_1^* \circ \overline{c} \longrightarrow j_2^*
\quad\text{and}\quad
j_2^* \circ \overline{d} \longrightarrow j_3^*
$$
of functors
$D^+_\QCoh(\mathcal{O}_{\overline{X}_2}) \to D^+_\QCoh(\mathcal{O}_X)$
and
$D^+_\QCoh(\mathcal{O}_{\overline{X}_3}) \to D^+_\QCoh(\mathcal{O}_X)$
for the same reasons as above. Denote $\overline{e}$ the
right adjoint of Lemma \ref{lemma-twisted-inverse-image}
for $h \circ g$. There is a canonical transformation
$$
j_1^* \circ \overline{e} \longrightarrow j_3^*
$$
of functors
$D^+_\QCoh(\mathcal{O}_{\overline{X}_3}) \to D^+_\QCoh(\mathcal{O}_X)$
given by (\ref{equation-sheafy}). Spelling things out we have to
show that the composition
$$
\alpha_h \circ \alpha_g :
j_1^* \circ \overline{a}_1 \to
j_1^* \circ \overline{c} \circ \overline{a}_2 \to
j_2^* \circ \overline{a}_2 \to
j_2^* \circ \overline{d} \circ \overline{a}_3 \to
j_3^* \circ \overline{a}_3
$$
is the same as the composition
$$
\alpha_{h \circ g} :
j_1^* \circ \overline{a}_1 \to
j_1^* \circ \overline{e} \circ \overline{a}_3 \to
j_3^* \circ \overline{a}_3
$$
We split this into two parts. The first is to show that the diagram
$$
\xymatrix{
\overline{a}_1 \ar[r] \ar[d] & \overline{c} \circ \overline{a}_2 \ar[d] \\
\overline{e} \circ \overline{a}_3 \ar[r] &
\overline{c} \circ \overline{d} \circ \overline{a}_3
}
$$
commutes where the lower horizontal arrow comes from the identification
$\overline{e} = \overline{c} \circ \overline{d}$. This is true
because the corresponding diagram of total direct image functors
$$
\xymatrix{
R\overline{f}_{1, *} \ar[r] \ar[d] & Rg_* \circ R\overline{f}_{2, *} \ar[d] \\
R(h \circ g)_* \circ R\overline{f}_{3, *} \ar[r] &
Rg_* \circ Rh_* \circ R\overline{f}_{3, *}
}
$$
is commutative (insert future reference here). The second part
is to show that the composition
$$
j_1^* \circ \overline{c} \circ \overline{d} \to
j_2^* \circ \overline{d} \to j_3^*
$$
is equal to the map
$$
j_1^* \circ \overline{e} \to j_3^*
$$
via the identification $\overline{e} = \overline{c} \circ \overline{d}$.
This was proven in Lemma \ref{lemma-compose-base-change-maps}
(note that in the current case the morphisms $f', g'$ of that
lemma are equal to $\text{id}_X$).
\end{proof}
